\section{偏微分方程}

偏微分方程把场在空间各点之间的联系与边界、初始条件同时编码。处理这类问题的主线有三条：先用自伴本征值问题建立正交函数系，再用分离变量把偏微分方程化成常微分方程，最后用正交展开或傅里叶变换确定各阶模态的系数。本节依次讨论 Sturm--Liouville 理论、三族正交多项式、波动方程与热方程，并给出所用定理的证明与例题的完整解答。

\subsection{Sturm--Liouville 问题与正交函数系}

先看固定在 $0<x<l$ 上的弦。分离变量后空间函数满足
\[
 y''+\lambda y=0,\qquad y(0)=y(l)=0.
\]
非零解只在 $\lambda_n=(n\pi/l)^2$ 时存在。取
\[
 y_n(x)=\sqrt{\frac2l}\sin\frac{n\pi x}{l},\qquad n=1,2,\ldots,
\]
便有
\[
 \int_0^l y_ny_m\,\dd x=\delta_{nm}.
\]
若 $y=\sum_n c_ny_n$，则 $c_n=\int_0^l yy_n\,\dd x$；在完备性成立时，Parseval 等式给出
$\int_0^l|y|^2\dd x=\sum_n|c_n|^2$。这正是有限维正交投影在函数空间中的延伸。

\begin{advdefinition}{带权内积、范数与完备正交系}\label{advA3:def:weighted-inner-product}
在区间 $(a,b)$ 上给定连续正权函数 $\rho(x)$，定义
\[
 (f,g)_\rho=\int_a^b f^*(x)g(x)\rho(x)\,\dd x,
 \qquad \|f\|_\rho=\sqrt{(f,f)_\rho}.
\]
若 $(z_i,z_j)_\rho=\delta_{ij}$，则 $\{z_i\}$ 是正交归一函数系；若每个属于所讨论函数空间的 $f$ 都能按范数收敛地写成 $f=\sum_i(z_i,f)_\rho z_i$，则称该函数系完备。
\end{advdefinition}

内积具有共轭对称性、线性和正定性。实函数情形可以去掉共轭号。完备性不是“每两个函数正交”自动带来的性质，它还要求这一组函数没有漏掉空间中的任何独立方向。

\begin{advexamplebox}{Cauchy--Schwarz 不等式与 Gram--Schmidt 正交化}\label{advA3:ex:cauchy-schmidt}
证明
\[
 |(f,g)_\rho|\le \|f\|_\rho\|g\|_\rho,
\]
并说明如何把线性无关函数 $z_1,\ldots,z_n$ 化为正交归一组。
\end{advexamplebox}
\begin{advsolution}
若 $g=0$ 结论显然。若 $g\ne0$，对任意复数 $c$ 有
\[
 0\le\|f-cg\|^2
 =\|f\|^2-c(f,g)-c^*(g,f)+|c|^2\|g\|^2.
\]
取 $c=(g,f)/\|g\|^2$，得到
\[
 0\le\|f\|^2-\frac{|(f,g)|^2}{\|g\|^2},
\]
即所证不等式；等号当且仅当 $f,cg$ 线性相关。

正交化递推为
\[
 u_1=z_1,\qquad y_1=\frac{u_1}{\|u_1\|},
\]
\[
 u_k=z_k-\sum_{j=1}^{k-1}(y_j,z_k)y_j,\qquad
 y_k=\frac{u_k}{\|u_k\|}.
\]
对 $i<k$，$(y_i,u_k)=(y_i,z_k)-(y_i,z_k)=0$。线性无关保证 $u_k\ne0$，因此归一化合法。每一步只是从 $z_k$ 中扣掉已知方向的投影，所以张成空间不变，最终 $\{y_k\}$ 是原空间的一组正交归一基。
\end{advsolution}

把 $1,x,x^2,\ldots$ 依次正交化可得到正交多项式。若 $Q_n$ 是 $n$ 次正交多项式，则它在 $(a,b)$ 内恰有 $n$ 个互异零点：若变号零点只有 $m<n$ 个，令 $R_m$ 为这些零点对应的一次因子之积，则 $Q_nR_m$ 在区间内保持同号，故 $(Q_n,R_m)_\rho\ne0$，与 $Q_n$ 和全部低于 $n$ 次多项式正交矛盾。常用三族为
\[
\begin{array}{c|c|c}
\text{多项式}&\text{区间}&\rho(x)\\\hline
\text{Legendre}&(-1,1)&1\\
\text{Hermite}&(-\infty,\infty)&\e^{-x^2}\\
\text{Laguerre}&(0,\infty)&x^\mu\e^{-x},\ \mu>-1
\end{array}
\]

\begin{advdefinition}{Sturm--Liouville 问题}\label{advA3:def:sl-problem}
给定区间 $[a,b]$ 上 $p(x)>0$、$\rho(x)>0$ 和实函数 $q(x)$，求
\[
 \frac{\dd}{\dd x}\!\left(p\frac{\dd y}{\dd x}\right)
 +q(x)y+\lambda\rho(x)y=0                                      \tag{A3.SL}
\]
在分离边界条件
\[
 c_1y(a)+c'_1y'(a)=0,\qquad c_2y(b)+c'_2y'(b)=0
\]
或周期边界条件 $y(a)=y(b),\ p(a)y'(a)=p(b)y'(b)$ 下的非零解。使非零解存在的 $\lambda$ 称本征值，相应的 $y$ 称本征函数。无穷区间还要求 $\int |y|^2\rho\dd x<\infty$。
\end{advdefinition}

定义算符
\[
 L=-\frac1\rho\left[\frac{\dd}{\dd x}
 \left(p\frac{\dd}{\dd x}\right)+q\right],
\]
则方程写成 $Ly=\lambda y$。固定弦、Legendre、Hermite 和 Laguerre 方程分别对应
\[
\begin{array}{c|c}
\text{问题}&\text{自伴形式}\\\hline
\text{固定弦}&y''+\lambda y=0\\
\text{Legendre}&[(1-x^2)y']'+\lambda y=0\\
\text{Hermite}&(\e^{-x^2}y')'+\lambda\e^{-x^2}y=0\\
\text{Laguerre}&(x^{\mu+1}\e^{-x}y')'+\lambda x^\mu\e^{-x}y=0
\end{array}
\]

\begin{advtheorem}{Sturm--Liouville 算符的自伴性}\label{advA3:thm:sl-self-adjoint}
在上述边界条件及允许的函数空间内，
\[
 (f,Lg)_\rho=(Lf,g)_\rho.
\]
\end{advtheorem}
\noindent\textbf{证明}\quad 两次分部积分等价于 Lagrange 恒等式
\[
 (f,Lg)_\rho-(Lf,g)_\rho
 =\bigl[p(f'^*g-f^*g')\bigr]_a^b.                         \tag{A3.1}
\]
分离边界条件在每个端点都把 $f,f'$ 与 $g,g'$ 限制为同一线性关系，因此括号为零；周期边界条件下两端贡献相消。奇异端点（例如 Legendre 的 $x=\pm1$）由 $p=0$ 消去边界项；无穷端点则由平方可积性和允许的衰减条件消去边界项，故式(A3.1)右端为零。\hfill$\square$

\begin{advtheorem}{正则 Sturm--Liouville 谱的存在性与完备性}\label{advA3:thm:sl-spectrum}
正则分离边界问题的本征值为离散实数列
\[
 \lambda_1<\lambda_2<\cdots<\lambda_n<\cdots,\qquad \lambda_n\to+\infty.
\]
每个本征值的本征空间为一维；可取实归一化本征函数 $\phi_n$，它在 $(a,b)$ 内恰有 $n-1$ 个零点。$\{\phi_n\}$ 在相应带权 $L^2$ 空间中完备。
\end{advtheorem}
这个谱定理保证后面的模态展开不是形式猜测。正则问题的一维本征空间来自二阶方程的 Wronskian 与两端边界条件；离散性和完备性可由自伴算符的 Green 函数是紧算符来证明。

\begin{advtheorem}{不同本征值的本征函数正交}\label{advA3:thm:sl-orthogonality}
归一化后
\[
 (\phi_n,\phi_m)_\rho=\delta_{nm}.
\]
\end{advtheorem}
\noindent\textbf{证明}\quad 由自伴性和 $L\phi_j=\lambda_j\phi_j$，
\[
 \lambda_m(\phi_n,\phi_m)=(\phi_n,L\phi_m)
 =(L\phi_n,\phi_m)=\lambda_n(\phi_n,\phi_m).
\]
若 $n\ne m$，则 $\lambda_n\ne\lambda_m$，内积只能为零；再逐个归一化即得。\hfill$\square$

\begin{advtheorem}{正交多项式的特征方程}\label{advA3:thm:orthogonal-polynomial-ode}
若权函数满足 Pearson 方程
\[
 \frac{\rho'}\rho=\frac{\alpha(x)}{\beta(x)}
 =\frac{\alpha_0+\alpha_1x}{\beta_0+\beta_1x+\beta_2x^2},
\]
且 $p=\rho\beta$ 在两端消失，则 $n$ 次正交多项式 $Q_n$ 满足
\[
 \frac{\dd}{\dd x}\left(pQ_n'\right)+\lambda_n\rho Q_n=0,
 \qquad \boxed{\lambda_n=-n(n+1)\beta_2-n\alpha_1}.
\]
\end{advtheorem}
\noindent\textbf{证明}\quad 对 $k<n$ 两次分部积分，边界项由 $p(a)=p(b)=0$ 消失：
\[
 \int_a^b (pQ_n')'x^k\dd x
 =k\int_a^bQ_n\{(\beta'+\alpha)x^{k-1}+(k-1)\beta x^{k-2}\}\rho\dd x=0.
\]
括号内至多为 $k$ 次多项式。因此 $(pQ_n')'/\rho$ 与所有低于 $n$ 次的多项式正交，它本身又至多是 $n$ 次，只能等于 $-\lambda_nQ_n$。比较 $x^n$ 的最高次项，得到
$\beta_2n(n-1)+(2\beta_2+\alpha_1)n+\lambda_n=0$，即所列本征值。\hfill$\square$

\begin{advexamplebox}{把一般二阶常微分方程化为自伴形式}\label{advA3:ex:ode-to-sl}
设
\[
 P(x)y''+Q(x)y'+[R_0(x)+\lambda R_1(x)]y=0,
 \]
其中 $P\ne0$。求等价的 Sturm--Liouville 形式。
\end{advexamplebox}
\begin{advsolution}
先除以 $P$，再乘积分因子 $\mu$。要使 $y''$ 与 $y'$ 合成 $(py')'$，须有
\[
 p=\mu,\qquad p'=\mu\frac QP,\qquad
 \frac{\mu'}\mu=\frac QP.
\]
故
\[
 \mu(x)=C\exp\!\left(\int^x\frac{Q(s)}{P(s)}\dd s\right).
\]
于是
\[
 (\mu y')'+\mu\frac{R_0}{P}y+\lambda\mu\frac{R_1}{P}y=0,
\]
即
\[
 \boxed{p=\mu,\quad q=\mu R_0/P,\quad \rho=\mu R_1/P}.
\]
若要成为正则 SL 问题，还需检查所选区间上 $p,\rho>0$ 以及边界条件确实消去边界型。各项量纲都等于 $\mu y''$ 的量纲。
\end{advsolution}

\begin{advexamplebox}{Bessel 方程的自伴形式}\label{advA3:ex:bessel-sl}
将
\[
 x^2y''+xy'+(x^2-\alpha^2)y=0
\]
化为自伴形式，并判断它是否存在一族有限次正交多项式解。
\end{advexamplebox}
\begin{advsolution}
除以 $x$ 即得
\[
 \boxed{(xy')'+\left(x-\frac{\alpha^2}{x}\right)y=0}.
\]
若把 $\alpha^2$ 当本征参数，可写成
$(xy')'+xy-\alpha^2x^{-1}y=0$；若把固定区间边界条件决定的径向波数 $k$ 当本征参数，则标准形式是
\[
 (xy')'+\left(k^2x-\frac{\alpha^2}{x}\right)y=0,
\]
权函数为 $\rho=x$。它不满足定理\ref{advA3:thm:orthogonal-polynomial-ode}所需的 $q=0$（或可并入常数本征值）条件，幂级数也不会在一般阶次截断，故解为 Bessel 函数而不是有限次正交多项式。只有特殊退化参数可能偶然化简，不能构成完整多项式族。
\end{advsolution}

\begin{advexamplebox}{Chebyshev 多项式}\label{advA3:ex:chebyshev}
讨论
\[
 (1-x^2)y''-xy'+\lambda y=0,\qquad -1<x<1
\]
在自然边界条件 $\lim_{x\to\pm1}\sqrt{1-x^2}\,y'=0$ 下的本征值，并写出零至四次的归一化多项式解。
\end{advexamplebox}
\begin{advsolution}
令 $x=\cos\theta$，则
\[
 (1-x^2)y''-xy'=\frac{\dd^2y}{\dd\theta^2}.
\]
边界条件转化为 $y_\theta(0)=y_\theta(\pi)=0$。对 $y_{\theta\theta}+\lambda y=0$ 求解可得 $\lambda\ge0$，且 $\sqrt\lambda$ 必须为整数，故
$y=\cos n\theta$、$\lambda_n=n^2$。因此第一类 Chebyshev 多项式
\[
 T_n(x)=\cos(n\arccos x)
\]
的前五个为
\[
 T_0=1, T_1=x, T_2=2x^2-1, T_3=4x^3-3x,
\ T_4=8x^4-8x^2+1.
\]
仅要求端点有限不足以确定这些本征值，例如 $\cos(\alpha\arccos x)$ 对非整数 $\alpha$ 也有限，但不满足上述两端条件。其权函数为 $\rho=(1-x^2)^{-1/2}$，且
\[
 \int_{-1}^1\frac{T_nT_m}{\sqrt{1-x^2}}\dd x
 =\begin{cases}\pi,&n=m=0,\\\pi/2,&n=m\ge1,\\0,&n\ne m.
 \end{cases}
\]
故归一化解是 $\widehat T_0=T_0/\sqrt\pi$、
$\widehat T_n=\sqrt{2/\pi}\,T_n$（$n\ge1$）。代回原方程可逐项检查；$\lambda_n$ 无量纲，与角变量二阶导数一致。
\end{advsolution}

\begin{advtheorem}{Rodrigues 公式}\label{advA3:thm:rodrigues}
在定理\ref{advA3:thm:orthogonal-polynomial-ode}的条件下，一组与 $Q_n$ 成比例的正交多项式为
\[
 \boxed{Q_n(x)=\frac1{\rho(x)}\frac{\dd^n}{\dd x^n}
 \bigl[\rho(x)\beta^n(x)\bigr]}.
\]
\end{advtheorem}
\noindent\textbf{证明}\quad 令
\[
 Q_n^{(k)}=\frac1{\rho\beta^{n-k}}
 \frac{\dd^k}{\dd x^k}(\rho\beta^n),\qquad k=0,1,\ldots,n.
\]
由 $\rho'\beta=\rho\alpha$，
\[
 Q_n^{(k+1)}=[\alpha+(n-k)\beta']Q_n^{(k)}
 +\beta(Q_n^{(k)})',
\]
所以其次数逐步增加，$Q_n^{(n)}$ 是 $n$ 次多项式。对 $m<n$ 分部积分 $n$ 次，端点项由 $\rho\beta^j$ 的衰减消失，
\[
 \int_a^b x^mQ_n\rho\dd x
 =(-1)^n\int_a^b\frac{\dd^n x^m}{\dd x^n}\rho\beta^n\dd x=0.
\]
故它与所有低次多项式正交。\hfill$\square$

\begin{advtheorem}{正交多项式的母函数}\label{advA3:thm:generating-function}
令 $z_1=z_1(z,\varepsilon)$ 是方程
\[
 \zeta=z+\varepsilon\beta(\zeta)
\]
中当 $\varepsilon\to0$ 时趋近 $z$ 的根，则 Rodrigues 多项式的指数型母函数为
\[
 \boxed{\Psi(z,\varepsilon)=
 \frac{\rho(z_1)}{\rho(z)}\frac1{1-\varepsilon\beta'(z_1)}
 =\sum_{n=0}^\infty\frac{Q_n(z)}{n!}\varepsilon^n}.
\]
\end{advtheorem}
\noindent\textbf{证明}\quad 把 Rodrigues 公式代入级数并用 Cauchy 积分公式，
\[
 \Psi=\frac1{2\pi\ii\rho(z)}\oint
 \frac{\rho(\zeta)}{\zeta-z-\varepsilon\beta(\zeta)}\dd\zeta.
\]
取只包围邻近根 $z_1$ 的小围道。分母对 $\zeta$ 的导数为
$1-\varepsilon\beta'(z_1)$，由留数定理立刻得到结论。这个证明同时说明级数的收敛范围由另一个根或 $\rho$ 的奇点决定。\hfill$\square$

\subsection{Legendre、Hermite 与 Laguerre 多项式}

对 Legendre 方程，$\rho=1$、$\beta=1-x^2$，故
\[
 \lambda_n=n(n+1),\qquad
 P_n(x)=\frac1{2^nn!}\frac{\dd^n}{\dd x^n}(x^2-1)^n,
\]
\[
 \frac1{\sqrt{1-2xt+t^2}}=\sum_{n=0}^\infty P_n(x)t^n.
\]
对母函数求导并比较 $t^n$ 系数，得到 Bonnet 递推式
\[
 (2n+1)xP_n=(n+1)P_{n+1}+nP_{n-1}.
\]
又由 $P_n(1)=1$、$P_n(-1)=(-1)^n$ 和正交性，
\[
 \int_{-1}^1P_nP_m\dd x=\frac2{2n+1}\delta_{nm}.
\]
前六个多项式为
\[
 P_0=1,\quad P_1=x,\quad P_2=\frac{3x^2-1}{2},\quad
 P_3=\frac{5x^3-3x}{2},
\]
\[
 P_4=\frac{35x^4-30x^2+3}{8},\qquad
 P_5=\frac{63x^5-70x^3+15x}{8}.
\]

\begin{center}\begin{minipage}{\linewidth}\centering
\centering
\begin{tikzpicture}
\begin{axis}[width=.82\textwidth,height=6.2cm,domain=-1:1,xmin=-1,xmax=1,ymin=-1.08,ymax=1.08,
 axis lines=middle,grid=major,xlabel={$x$},ylabel={$P_n(x)$},
 legend style={font=\scriptsize,at={(0.5,-0.18)},anchor=north,legend columns=3},
 every axis plot/.append style={line width=.8pt},samples=181]
 \addplot[DeepBlue] {1}; \addlegendentry{$P_0$}
 \addplot[orange!80!black] {x}; \addlegendentry{$P_1$}
 \addplot[green!55!black] {(3*x^2-1)/2}; \addlegendentry{$P_2$}
 \addplot[red!75!black] {(5*x^3-3*x)/2}; \addlegendentry{$P_3$}
 \addplot[violet] {(35*x^4-30*x^2+3)/8}; \addlegendentry{$P_4$}
 \addplot[brown!80!black] {(63*x^5-70*x^3+15*x)/8}; \addlegendentry{$P_5$}
\end{axis}
\end{tikzpicture}
\bookcaption{前六个 Legendre 多项式。每个 $P_n$ 在 $(-1,1)$ 内恰有 $n$ 个零点，端点值固定为 $\pm1$。}\label{advA3:fig:legendre}
\end{minipage}\end{center}

Hermite 方程给出
\[
 \lambda_n=2n,\qquad H_n(x)=(-1)^n\e^{x^2}\frac{\dd^n}{\dd x^n}\e^{-x^2},
\]
\[
 \e^{2xt-t^2}=\sum_{n=0}^\infty H_n(x)\frac{t^n}{n!},\qquad
 H_n'=2nH_{n-1},\qquad 2xH_n=H_{n+1}+2nH_{n-1}.
\]
利用 Rodrigues 公式分部积分 $n$ 次，
\[
 \int_{-\infty}^{\infty}H_nH_m\e^{-x^2}\dd x
 =2^nn!\sqrt\pi\,\delta_{nm}.
\]
前六个为
$H_0=1$、$H_1=2x$、$H_2=4x^2-2$、$H_3=8x^3-12x$、
$H_4=16x^4-48x^2+12$、$H_5=32x^5-160x^3+120x$。

\begin{center}\begin{minipage}{\linewidth}\centering
\centering
\begin{tikzpicture}
\begin{axis}[width=.82\textwidth,height=6.2cm,domain=-2.1:2.1,xmin=-2.1,xmax=2.1,ymin=-45,ymax=50,
 axis lines=middle,grid=major,xlabel={$x$},ylabel={$H_n(x)$},restrict y to domain=-55:60,
 legend style={font=\scriptsize,at={(0.5,-0.18)},anchor=north,legend columns=3},samples=221]
 \addplot[DeepBlue] {1}; \addlegendentry{$H_0$}
 \addplot[orange!80!black] {2*x}; \addlegendentry{$H_1$}
 \addplot[green!55!black] {4*x^2-2}; \addlegendentry{$H_2$}
 \addplot[red!75!black] {8*x^3-12*x}; \addlegendentry{$H_3$}
 \addplot[violet] {16*x^4-48*x^2+12}; \addlegendentry{$H_4$}
 \addplot[brown!80!black] {32*x^5-160*x^3+120*x}; \addlegendentry{$H_5$}
\end{axis}
\end{tikzpicture}
\bookcaption{前六个物理学家约定的 Hermite 多项式；纵轴截断仅为同时显示低阶曲线。}\label{advA3:fig:hermite}
\end{minipage}\end{center}

$\mu$ 阶 Laguerre 方程给出
\[
 xL_{n,\mu}''+(\mu+1-x)L_{n,\mu}'+nL_{n,\mu}=0,
\]
\[
 L_{n,\mu}(x)=\frac{x^{-\mu}\e^x}{n!}
 \frac{\dd^n}{\dd x^n}(\e^{-x}x^{n+\mu}),\qquad
 \frac{\e^{-xt/(1-t)}}{(1-t)^{\mu+1}}
 =\sum_{n=0}^\infty L_{n,\mu}(x)t^n.
\]
其范数和递推关系为
\[
 \int_0^\infty L_{n,\mu}L_{m,\mu}x^\mu\e^{-x}\dd x
 =\frac{\Gamma(n+\mu+1)}{n!}\delta_{nm},
\]
\[
 xL_{n,\mu}=(2n+\mu+1)L_{n,\mu}
 -(n+\mu)L_{n-1,\mu}-(n+1)L_{n+1,\mu},
\]
\[
 \frac{\dd^k}{\dd x^k}L_{n,\mu}=(-1)^kL_{n-k,\mu+k}\quad(k\le n).
\]
前四个为
\[
 L_{0,\mu}=1,\qquad L_{1,\mu}=-x+\mu+1,
\]
\[
 L_{2,\mu}=\frac{x^2}{2}-(\mu+2)x+\frac{(\mu+1)(\mu+2)}2,
\]
\[
 L_{3,\mu}=-\frac{x^3}{6}+\frac{(\mu+3)x^2}{2}
 -\frac{(\mu+2)(\mu+3)x}{2}
 +\frac{(\mu+1)(\mu+2)(\mu+3)}6.
\]

\begin{center}\begin{minipage}{\linewidth}\centering
\centering
\begin{tikzpicture}
\begin{axis}[width=.82\textwidth,height=6.2cm,domain=0:8,xmin=0,xmax=8,ymin=-5,ymax=8,
 axis lines=middle,grid=major,xlabel={$x$},ylabel={$L_{n,\mu}(x)$},restrict y to domain=-6:9,
 legend style={font=\scriptsize,at={(0.5,-0.18)},anchor=north,legend columns=4},samples=221]
 \addplot[DeepBlue] {1}; \addlegendentry{$L_{0,0}$}
 \addplot[orange!80!black] {1-x}; \addlegendentry{$L_{1,0}$}
 \addplot[green!55!black] {x^2/2-2*x+1}; \addlegendentry{$L_{2,0}$}
 \addplot[red!75!black] {-x^3/6+1.5*x^2-3*x+1}; \addlegendentry{$L_{3,0}$}
 \addplot[violet,dashed] {2-x}; \addlegendentry{$L_{1,1}$}
 \addplot[brown!80!black,dashed] {x^2/2-3*x+3}; \addlegendentry{$L_{2,1}$}
 \addplot[cyan!60!black,dashed] {-x^3/6+2*x^2-6*x+4}; \addlegendentry{$L_{3,1}$}
\end{axis}
\end{tikzpicture}
\bookcaption{若干广义 Laguerre 多项式。实线为 $\mu=0$，虚线为 $\mu=1$；每条曲线的次数和零点数可直接核对。}\label{advA3:fig:laguerre}
\end{minipage}\end{center}

\begin{advexamplebox}{Legendre 多项式的导数递推}\label{advA3:ex:legendre-derivative-recurrence}
证明
\[
 P'_{n+1}(x)-xP'_n(x)=(n+1)P_n(x).
\]
\end{advexamplebox}
\begin{advsolution}
从 Bonnet 递推式
\[
 (n+1)P_{n+1}=(2n+1)xP_n-nP_{n-1}
\]
求导，再使用上一阶关系消去 $P'_{n-1}$ 也可证明。更直接地，对母函数
$G=(1-2xt+t^2)^{-1/2}$ 求 $x$ 导数：
\[
 \partial_xG=t(1-2xt+t^2)^{-3/2}.
\]
另一方面
\[
 (1-xt)\partial_xG-t^2\partial_tG=tG.
\]
比较 $t^{n+1}$ 系数得到
$P'_{n+1}-xP'_n=(n+1)P_n$。在 $n=0$ 时左边是 $P'_1=1$，右边也是1，起始阶检查通过。
\end{advsolution}

\begin{advexamplebox}{Legendre 多项式导数的范数}\label{advA3:ex:legendre-derivative-norm}
证明
\[
 \int_{-1}^1[P'_n(x)]^2\dd x=n(n+1).
\]
\end{advexamplebox}
\begin{advsolution}
由导数递推反复展开，
\[
 P'_n=(2n-1)P_{n-1}+(2n-5)P_{n-3}+\cdots.
\]
各项彼此正交，故
\[
 \int_{-1}^1(P'_n)^2\dd x
 =\sum_{j\ge0}(2n-4j-1)^2\frac{2}{2(n-2j-1)+1}.
\]
每项约去一个因子，成为
$2\sum_{j\ge0}(2n-4j-1)$，末项取到正数1或3。这个等差和等于 $n(n+1)$。

也可用归纳：由 $P'_{n+1}=xP'_n+(n+1)P_n$，平方积分。分部积分给出
$\int xP'_nP_n\dd x=1-\tfrac12\int P_n^2\dd x$，再代入
$\int P_n^2\dd x=2/(2n+1)$ 和由微分方程导出的
$\int x^2(P'_n)^2=\int(P'_n)^2-n(n+1)\int P_n^2$，即可由 $n$ 推到 $n+1$。结果为纯数，与 $x$ 无量纲一致。
\end{advsolution}

\begin{advexamplebox}{Legendre 多项式的偶极矩}\label{advA3:ex:legendre-dipole}
计算
\[
 \mathcal P_{ij}=\int_{-1}^1xP_i(x)P_j(x)\dd x.
\]
\end{advexamplebox}
\begin{advsolution}
由 Bonnet 递推
\[
 xP_i=\frac{i+1}{2i+1}P_{i+1}+\frac{i}{2i+1}P_{i-1}.
\]
与 $P_j$ 内积并用正交性，只有 $j=i\pm1$ 非零：
\[
 \boxed{\mathcal P_{ij}=
 \frac{2(i+1)}{(2i+1)(2i+3)}\delta_{j,i+1}
 +\frac{2i}{(2i+1)(2i-1)}\delta_{j,i-1}}.
\]
交换 $i,j$ 后数值不变，符合积分本身的对称性；$i=j$ 时被积函数奇偶性也立刻给出零。
\end{advsolution}

三族多项式的核心数据可压缩为下表，后续量子力学和势函数展开会反复使用：
\begingroup\setlength{\arraycolsep}{3pt}
\[
\begin{array}{c|c|c|c}
&\text{Legendre}&\text{Hermite}&\text{Laguerre }(\mu=0)\\\hline
\text{方程}&(1-x^2)y''-2xy'+n(n+1)y=0&y''-2xy'+2ny=0&xy''+(1-x)y'+ny=0\\
\text{权函数}&1&\e^{-x^2}&\e^{-x}\\
\text{范数平方}&2/(2n+1)&2^nn!\sqrt\pi&1
\end{array}
\]
\endgroup

\subsection{常见二阶偏微分方程及其分类}

物理中最常见的线性二阶方程包括
\[
 \text{波动方程：}\quad \partial_t^2u-c^2\nabla^2u=0,
\]
\[
 \text{热方程：}\quad \partial_tu-\kappa\nabla^2u=0,
\]
\[
 \text{Schr\"odinger 方程：}\quad
 \ii\hbar\partial_t\psi=H(\vect r,-\ii\hbar\nabla)\psi,
\]
\[
 \text{Laplace/Poisson/Helmholtz 方程：}\quad
 \nabla^2\phi=0,\quad \nabla^2\phi=-\rho/\varepsilon_0,\quad
 \nabla^2\phi+k^2\phi=0.
\]
柱坐标与球坐标下的 Laplace 算符分别为
\[
 \nabla^2u=\frac1r\partial_r(r\partial_ru)
 +\frac1{r^2}\partial_\varphi^2u+\partial_z^2u,
\]
\[
 \nabla^2u=\frac1{r^2}\partial_r(r^2\partial_ru)
 +\frac1{r^2\sin\theta}\partial_\theta(\sin\theta\,\partial_\theta u)
 +\frac1{r^2\sin^2\theta}\partial_\varphi^2u.
\]
这些方程都是线性的，最高偏导数为二阶；除 Poisson 方程外，上列方程在右端为零时齐次。坐标变换会让常系数方程出现变系数，但不会改变其几何类型。

\begin{advdefinition}{二阶线性偏微分方程的类型}\label{advA3:def:pde-classification}
在一点把二阶主部写成
\[
 \sum_{i,j=1}^n a_{ij}\partial_i\partial_j u,
\]
其中 $A=(a_{ij})$ 取对称部分。若 $A$ 非退化且全部特征值同号，称为椭圆型；若 $A$ 有且只有一个零特征值、其余同号，称为抛物型；若 $A$ 非退化，且一正其余负或一负其余正，称为双曲型；若正、负特征值各多于一个，则称为超双曲型。更高维中若零特征值多于一个，还需视作退化型单独说明。
\end{advdefinition}

\begin{advexamplebox}{典型物理方程的分类}\label{advA3:ex:pde-types}
判断波动方程、热方程、Schr\"odinger 方程、Laplace 方程、Poisson 方程和 Helmholtz 方程的类型。
\end{advexamplebox}
\begin{advsolution}
把时间也视为自变量。$n$ 维波动方程主部矩阵为
$\operatorname{diag}(1,-c^2,\ldots,-c^2)$，有一正 $n$ 负，故为双曲型。热方程关于 $(t,\vect x)$ 的二阶主部是
$\operatorname{diag}(0,-\kappa,\ldots,-\kappa)$，恰有一个零特征值，故为抛物型。

自由 Schr\"odinger 方程具有虚数时间系数，不能直接套用上述实系数标量分类；应单列为色散型的 Schr\"odinger 方程。虽然对时间一阶、对空间二阶，但它不是热方程式的抛物演化：空间傅里叶模随时间乘以模为1的相位，而不是热方程的衰减因子。Laplace 方程的主部矩阵为单位阵，故为椭圆型；Poisson 的非齐次项不改变主部，仍为椭圆型；Helmholtz 的 $k^2\phi$ 是零阶项，也不改变类型，仍为椭圆型。分类只看最高阶项，因此势能、源项和常数质量项不改变结论。
\end{advsolution}

\subsection{分离变量、傅里叶分析与初边值问题}

\paragraph{先检查边界条件是否允许解存在。}
对 $-\Delta u=f$、外法向导数 $\partial_nu=g$，散度定理要求
\[
 \int_\Omega f\,\dd V+\oint_{\partial\Omega}g\,\dd S=0.
\]
纯 Neumann 问题若不满足此式便无解；满足时仍可加任意常数，需指定平均值来唯一化。它是守恒条件，不是计算结束后的可选检查。
例如 $-u''=f(x)$、$u'(0)=u'(L)=0$，必须有 $\int_0^L f\dd x=0$。若 $f=\cos(\pi x/L)$，零平均解为
\[
 u=\frac{L^2}{\pi^2}\cos\frac{\pi x}{L};
\]
若 $f=1$ 则不可能满足两端零通量。谱展开中同一条件表现为零本征值对应的常数模必须与源正交，不能把其系数除以零。


对一维波动方程
\[
 \partial_t^2y-c^2\partial_x^2y=0
\]
令 $y=f(x)g(t)$，则
\[
 \frac{g''}{g}=c^2\frac{f''}{f}=-\omega^2,\qquad k=\omega/c.
\]
于是
\[
 g=A\cos\omega t+B\sin\omega t,\qquad
 f=C\cos kx+D\sin kx.
\]
复指数形式把两列行波清楚分开：
\[
 y_+=C_+\e^{\ii(kx-\omega t)},\qquad
 y_-=C_-\e^{\ii(-kx-\omega t)},\qquad \omega=c|k|.
\]
若分离常数取相反符号，会出现 $\cosh$ 或 $\e^{\pm\beta x}$。在无限空间中这类解通常因不平方可积而舍去；在有限区域或倏逝波问题中则可能正是物理解。

\begin{advtheorem}{Fourier 反演与 Dirac $\delta$ 分布}\label{advA3:thm:fourier-inversion}
采用约定
\[
 \widetilde f(k)=\frac1{\sqrt{2\pi}}\int_{-\infty}^{\infty}f(x)\e^{-\ii kx}\dd x,\qquad
 f(x)=\frac1{\sqrt{2\pi}}\int_{-\infty}^{\infty}\widetilde f(k)\e^{\ii kx}\dd k,
\]
则分布意义下
\[
 \int_{-\infty}^{\infty}\e^{\ii kx}\dd x=2\pi\delta(k),
\]
其中 $\int\delta(k)\varphi(k)\dd k=\varphi(0)$。
\end{advtheorem}
\noindent\textbf{证明}\quad 用 Gaussian 核正则化：
\[
 I_\varepsilon(k)=\int_{-\infty}^{\infty}
 \e^{-\varepsilon x^2+\ii kx}\dd x
 =\sqrt{\frac\pi\varepsilon}\e^{-k^2/(4\varepsilon)}.
\]
对任意光滑快速衰减的 $\varphi$，令 $k=2\sqrt\varepsilon,u$，
\[
 \frac1{2\pi}\int I_\varepsilon(k)\varphi(k)\dd k
 =\frac1{\sqrt\pi}\int\e^{-u^2}\varphi(2\sqrt\varepsilon u)\dd u
 \longrightarrow\varphi(0).
\]
故 $I_\varepsilon\to2\pi\delta$，再代入 Fourier 变换的两重积分即得反演。\hfill$\square$

\begin{advexamplebox}{Dirichlet 积分}\label{advA3:ex:dirichlet-integral}
证明
\[
 \int_{-\infty}^{\infty}\frac{\sin kx}{x}\dd x=\pi\operatorname{sgn}k.
\]
\end{advexamplebox}
\begin{advsolution}
设 $k>0$，先加收敛因子
\[
 F_\varepsilon(k)=\int_0^\infty\e^{-\varepsilon x}
 \frac{\sin kx}{x}\dd x.
\]
对 $k$ 求导可以消去 $1/x$：
\[
 F_\varepsilon'(k)=\int_0^\infty\e^{-\varepsilon x}\cos kx\dd x
 =\frac{\varepsilon}{\varepsilon^2+k^2}.
\]
由 $F_\varepsilon(0)=0$，
$F_\varepsilon(k)=\arctan(k/\varepsilon)\to\pi/2$。被积函数为偶函数，因此全轴积分为 $\pi$。$k<0$ 时对 $k$ 为奇函数，结果反号；$k=0$ 时积分为零。由此得到统一表达式 $\pi\operatorname{sgn}k$。
\end{advsolution}

\begin{advtheorem}{d'Alembert 公式}\label{advA3:thm:dalembert}
无界弦初值问题
\[
 y_{tt}-c^2y_{xx}=0,\qquad y(x,0)=f(x),\quad y_t(x,0)=g(x)
\]
的解为
\[
 \boxed{y(x,t)=\frac{f(x-ct)+f(x+ct)}2
 +\frac1{2c}\int_{x-ct}^{x+ct}g(\xi)\dd\xi}.
\]
\end{advtheorem}
\noindent\textbf{证明}\quad 改用特征变量 $\xi=x-ct$、$\eta=x+ct$，则
\[
 \partial_t^2-c^2\partial_x^2=-4c^2\partial_\xi\partial_\eta.
\]
方程变成 $y_{\xi\eta}=0$，积分得 $y=F(\xi)+G(\eta)$。初始条件给出
\[
 F(x)+G(x)=f(x),\qquad -cF'(x)+cG'(x)=g(x).
\]
把第二式积分并与第一式联立，常数在 $F,G$ 中相消，得到所列公式。直接代入 $t=0$ 得 $y=f$；对 $t$ 求导再令 $t=0$ 得 $y_t=g$，方向和端点符号均得到核验。\hfill$\square$

有界弦的边界条件决定允许的空间本征函数。例如
\[
\begin{array}{c|c}
y(0)=y(l)=0&\sin(n\pi x/l),\ n\ge1\\
y(0)=0,\ y'(l)=0&\sin[(n+\tfrac12)\pi x/l],\ n\ge0\\
y'(0)=0,\ y(l)=0&\cos[(n+\tfrac12)\pi x/l],\ n\ge0\\
y'(0)=y'(l)=0&\cos(n\pi x/l),\ n\ge0
\end{array}
\]
若 $y(0,t)=0$、$y_x(l,t)=0$，则 $k_nl=(n+\tfrac12)\pi$，展开
\[
 y=\sum_{n=0}^\infty[A_n\cos(k_nct)+B_n\sin(k_nct)]\sin k_nx,
\]
其中
\[
 A_n=\frac2l\int_0^lf(x)\sin k_nx\dd x,\qquad
 B_n=\frac2{k_ncl}\int_0^lg(x)\sin k_nx\dd x.
\]
周期边界则同时需要 $\sin(2n\pi x/l)$ 与 $\cos(2n\pi x/l)$ 两支。

\begin{advdefinition}{$n$ 维 Fourier 变换}\label{advA3:def:nd-fourier}
对 $\mathbb R^n$ 上平方可积函数，采用对称归一化
\[
 \widetilde f(\vect k)=\frac1{(2\pi)^{n/2}}
 \int_{\mathbb R^n}f(\vect x)\e^{-\ii\vect k\cdot\vect x}\dd^n x,
\]
\[
 f(\vect x)=\frac1{(2\pi)^{n/2}}
 \int_{\mathbb R^n}\widetilde f(\vect k)\e^{\ii\vect k\cdot\vect x}\dd^n k.
\]
在这一约定下 $\mathcal F^{-1}=\mathcal F^*$，并有
$\mathcal F^2f(\vect x)=f(-\vect x)$。
\end{advdefinition}

\begin{advexamplebox}{Parseval 与 Plancherel 定理}\label{advA3:ex:parseval-plancherel}
证明
\[
 \int_{-\infty}^{\infty}f^*(x)g(x)\dd x
 =\int_{-\infty}^{\infty}\widetilde f^*(k)\widetilde g(k)\dd k,
\]
并由此得到 $\|f\|_2=\|\widetilde f\|_2$。
\end{advexamplebox}
\begin{advsolution}
把 $g$ 的反变换代入左边并交换积分次序：
\[
 \int f^*(x)g(x)\dd x
 =\frac1{\sqrt{2\pi}}\int\!\int
 f^*(x)\widetilde g(k)\e^{\ii kx}\dd k\dd x.
\]
括号中的 $x$ 积分是
\[
 \frac1{\sqrt{2\pi}}\int f^*(x)\e^{\ii kx}\dd x
 =\widetilde f^*(k),
\]
所以得到 Parseval 等式。令 $g=f$ 即得
$\int|f|^2=\int|\widetilde f|^2$，这就是 Plancherel 定理。它说明 Fourier 变换是 $L^2$ 空间的酉变换，量纲上两边还通过 $\widetilde f$ 所含的长度因子严格匹配。
\end{advsolution}

\begin{advexamplebox}{标准正态分布的 Fourier 变换}\label{advA3:ex:gaussian-fourier}
求
\[
 f(x)=\frac1{\sqrt{2\pi}}\e^{-x^2/2}
\]
在本节对称归一化约定下的 Fourier 变换。
\end{advexamplebox}
\begin{advsolution}
配方
\[
 -\frac{x^2}{2}-\ii kx=-\frac{(x+\ii k)^2}{2}-\frac{k^2}{2}.
\]
整函数 $\e^{-z^2/2}$ 的积分路径可平移，或直接使用 Gaussian 积分，因此
\[
 \widetilde f(k)=\frac1{2\pi}\int_{-\infty}^{\infty}
 \e^{-x^2/2-\ii kx}\dd x
 =\boxed{\frac1{\sqrt{2\pi}}\e^{-k^2/2}}.
\]
标准 Gaussian 在这一归一化下是 Fourier 变换的不动函数。取 $k=0$ 得
$\widetilde f(0)=1/\sqrt{2\pi}$，也等于 $(2\pi)^{-1/2}\int f\dd x$，常数检查一致。
\end{advsolution}

\begin{advexamplebox}{非齐次边界条件的齐次化}\label{advA3:ex:inhomogeneous-boundary}
讨论一维波动方程
\[
 y_{tt}-c^2y_{xx}=0,\qquad y(0,t)=a(t),\quad y(l,t)=b(t)
\]
的求解方法。
\end{advexamplebox}
\begin{advsolution}
在固定边界 $x=l$ 上，边值只能随时间变化。构造满足两端边界条件的线性函数
\[
 h(x,t)=\left(1-\frac xl\right)a(t)+\frac xl b(t),
\qquad u=y-h.
\]
于是 $u(0,t)=u(l,t)=0$，而
\[
 u_{tt}-c^2u_{xx}=-h_{tt}equiv F(x,t).
\]
展开
\[
 u(x,t)=\sum_{n=1}^\infty q_n(t)\sin\frac{n\pi x}{l},\qquad
 F_n(t)=\frac2l\int_0^lF(x,t)\sin\frac{n\pi x}{l}\dd x,
\]
得到每个模态
\[
 q_n''+\omega_n^2q_n=F_n(t),\qquad \omega_n=cn\pi/l.
\]
给定初值后用受迫振子公式求 $q_n$。最后 $y=u+h$。这样边界条件被精确满足，且 $a=b=0$ 时退回齐次边界结果。
\end{advsolution}

\begin{advexamplebox}{非齐次波动方程的模态解}\label{advA3:ex:forced-wave}
求解思路：
\[
 y_{tt}-c^2y_{xx}=A(x,t),\qquad y(0,t)=y(l,t)=0,
\]
并给出零初值时的显式模态公式。
\end{advexamplebox}
\begin{advsolution}
令
\[
 y=\sum_{n=1}^\infty q_n(t)\sin k_nx,\qquad
 A=\sum_{n=1}^\infty A_n(t)\sin k_nx,\qquad k_n=n\pi/l,
\]
其中
$A_n(t)=\tfrac2l\int_0^lA(x,t)\sin k_nx\dd x$。正交投影后
\[
 q_n''+\omega_n^2q_n=A_n(t),\qquad \omega_n=ck_n.
\]
零初值的 Green 函数解为
\[
 \boxed{q_n(t)=\frac1{\omega_n}\int_0^t
 A_n(\tau)\sin[\omega_n(t-\tau)]\dd\tau}.
\]
所以
\[
 y(x,t)=\sum_{n=1}^\infty\frac{\sin k_nx}{\omega_n}
 \int_0^tA_n(\tau)\sin[\omega_n(t-\tau)]\dd\tau.
\]
逐项代回时，每个模态都产生 $A_n(t)$；在 $t=0$ 时积分及其一阶导数均为零，初边值条件全部满足。
\end{advsolution}

\begin{advexamplebox}{无界域热方程的热核}\label{advA3:ex:heat-kernel}
证明初值问题
\[
 u_t-\kappa u_{xx}=0,\qquad u(x,0)=f(x),\quad x\in\mathbb R
\]
的解为
\[
 u(x,t)=\frac1{\sqrt{4\pi\kappa t}}
 \int_{-\infty}^{\infty}
 \exp\!\left[-\frac{(x-\xi)^2}{4\kappa t}\right]f(\xi)\dd\xi.
\]
\end{advexamplebox}
\begin{advsolution}
对 $x$ 作 Fourier 变换，$\widetilde{u_{xx}}=-k^2\widetilde u$，故
\[
 \partial_t\widetilde u+\kappa k^2\widetilde u=0,\qquad
 \widetilde u(k,t)=\widetilde f(k)\e^{-\kappa k^2t}.
\]
反变换中乘积对应卷积，而
\[
 \frac1{2\pi}\int_{-\infty}^{\infty}
 \e^{-\kappa k^2t+\ii kx}\dd k
 =\frac1{\sqrt{4\pi\kappa t}}\e^{-x^2/(4\kappa t)}
 \equiv K_t(x).
\]
于是 $u=K_t*f$，即题给公式。核满足
$\int K_t\dd x=1$，所以总量 $\int u\dd x$ 守恒；$\kappa t$ 的量纲是长度平方，指数无量纲。$t\downarrow0$ 时 $K_t\to\delta$，故恢复初值。直接计算也有 $\partial_tK_t=\kappa\partial_x^2K_t$，方程校验完成。
\end{advsolution}

\subsection{球坐标分离与球谐函数}

对球坐标中的 Laplace 方程令 $\phi=R(r)Y(\theta,\varphi)$，得到
\[
 \frac1R\frac{\dd}{\dd r}\left(r^2\frac{\dd R}{\dd r}\right)
 =-\frac1Y\left[
 \frac1{\sin\theta}\partial_\theta(\sin\theta\,\partial_\theta Y)
 +\frac1{\sin^2\theta}\partial_\varphi^2Y\right]
 =l(l+1).
\]
径向解为
\[
 R(r)=A_lr^l+B_lr^{-l-1}.
\]
若轴对称，角方程在 $x=\cos\theta$ 后成为 Legendre 方程，要求两极有限迫使 $l$ 为非负整数，故
\[
 \phi(r,\theta)=\sum_{l=0}^\infty
 (A_lr^l+B_lr^{-l-1})P_l(\cos\theta).
\]
若还有 $\varphi$ 依赖，再令 $Y=y(x)z(\varphi)$。周期性给出
$z=\e^{\ii m\varphi}$、$m\in\mathbb Z$；$y$ 满足连带 Legendre 方程
\[
 \frac{\dd}{\dd x}\left[(1-x^2)y'\right]
 +\left[l(l+1)-\frac{m^2}{1-x^2}\right]y=0.
\]
其有限解可写为
\[
 P_l^m(x)=(-1)^m(1-x^2)^{m/2}\frac{\dd^mP_l}{\dd x^m},
 \qquad |m|\le l.
\]
这里采用含 Condon--Shortley 相位 $(-1)^m$ 的现代约定；若去掉该相位，下面负阶公式中的相应符号也要同步改变。

\begin{advexamplebox}{连带 Legendre 函数的负阶与宇称关系}\label{advA3:ex:associated-legendre}
证明
\[
 P_l^{-m}(x)=(-1)^m\frac{(l-m)!}{(l+m)!}P_l^m(x),
 \qquad P_l^m(-x)=(-1)^{l+m}P_l^m(x).
\]
\end{advexamplebox}
\begin{advsolution}
第二式最直接。由 Rodrigues 公式 $P_l(-x)=(-1)^lP_l(x)$，对 $x$ 求 $m$ 阶导数时每次链式法则再带一个负号，故
\[
 \frac{\dd^mP_l(-x)}{\dd(-x)^m}=(-1)^{l+m}
 \frac{\dd^mP_l(x)}{\dd x^m}.
\]
乘上偶函数 $(1-x^2)^{m/2}$，得到
$P_l^m(-x)=(-1)^{l+m}P_l^m(x)$。

对负阶，连带 Legendre 方程只含 $m^2$，所以 $P_l^{-m}$ 与 $P_l^m$ 是同一端点正则解的常数倍。用 Rodrigues 公式比较 $x\to1^-$ 的首项，或比较球谐函数归一化
\[
 Y_l^m=N_{lm}P_l^m(\cos\theta)\e^{\ii m\varphi},\qquad
 Y_l^{-m}=(-1)^m(Y_l^m)^*,
\]
可得比例常数
\[
 \boxed{P_l^{-m}=(-1)^m\frac{(l-m)!}{(l+m)!}P_l^m}.
\]
例如 $l=1,m=1$ 时 $P_1^1=-\sqrt{1-x^2}$、
$P_1^{-1}=\tfrac12\sqrt{1-x^2}$，正好满足公式，符号约定检查通过。
\end{advsolution}
